% Lösungen Einheit 1 von 5 – Punkte darstellen und Lage lesen (zusammenbau v0.1)
\einheitenkopf{Einheit 1 von 5 · Punkte darstellen und Lage lesen}

\erg{11}{a) $P$ liegt in der $x_1x_2$-Ebene \quad b) $C(4 | 5 | 2)$; $A$ und $B$ siehe Lösungsgrafik \quad c) Jeder Punkt hat zwei Koordinaten $0$ oder $4$ und liegt damit auf einer Würfelkante; die vier Punkte der Reihe nach verbinden, siehe Lösungsgrafik. \quad d) $A$ und $D$ haben eine negative $x_2$-Koordinate und liegen links der $x_1x_3$-Ebene, $C$ und $D$ haben eine negative $x_1$-Koordinate und liegen hinter der $x_2x_3$-Ebene; siehe Lösungsgrafik. \quad e) $P$ liegt auf der Deckkante mit $x_1 = 5$, ein Drittel ihrer Länge von vorn; Dreieck siehe Lösungsgrafik. \quad f) Die dritten Koordinaten $3$ und $2$ sind beide positiv: beide Punkte liegen oberhalb der $x_1x_2$-Ebene. \quad g) Fehlende Ecken $(-5 | 3)$, $(-6 | 0)$, $(-5 | -3)$, $(0 | -4)$, $(5 | -3)$; siehe Lösungsgrafik. \quad h) $|SB| = |SD| = \sqrt{9 + 16} = 5$ cm (nicht $4$ cm); Spitzen bei $(-4 | 0)$ an $AD$, $(8 | 0)$ an $BC$ und $(0 | 8)$ an $CD$ – die Dreiecke an $BC$ und $CD$ sind bei $B$ bzw. $D$ rechtwinklig. \quad i) $B'(5 | 1 | 0)$, $C'(1 | 4 | 0)$, $M'(2 | 2 | 0)$; $A'$ und $D'$ durch Spiegeln an $M'$ (nicht durch Verschieben): $A'(3 | 0 | 0)$, $D'(-1 | 3 | 0)$.}

11b)

\begin{ksys3} \rpunkt{5}{3}{1}{A} \rpunkt{2}{1}{4}{B} \rpunkt{4}{5}{2}{C} \end{ksys3}


11c)

\begin{ksys3} \rquader{0,0,0}{4}{4}{4} \rpunkt*{4}{0}{1}{P} \rpunkt*{1}{4}{0}{Q} \rpunkt*{0}{4}{2}{R} \rpunkt*{3}{0}{4}{S} \rgerade[0:1]{4,0,1}{-3,4,-1}{} \rgerade[0:1]{1,4,0}{-1,0,2}{} \rgerade[0:1]{0,4,2}{3,-4,2}{} \rgerade[0:1]{3,0,4}{1,0,-3}{} \end{ksys3}


11d)

\begin{ksys3} \rpunkt{5}{-3}{0}{A} \rpunkt{5}{2}{0}{B} \rpunkt{-1}{4}{4}{C} \rpunkt{-1}{-1}{4}{D} \rgerade[0:1]{5,-3,0}{0,5,0}{} \rgerade[0:1]{5,2,0}{-6,2,4}{} \rgerade[0:1]{-1,4,4}{0,-5,0}{} \rgerade[0:1]{-1,-1,4}{6,-2,-4}{} \end{ksys3}


11e)

\begin{ksys3} \rquader{0,0,0}{5}{6}{3} \rpunkt*{0}{0}{0}{O} \rpunkt*{5}{2}{3}{P} \rpunkt*{0}{6}{3}{R} \rgerade[0:1]{0,0,0}{5,2,3}{} \rgerade[0:1]{5,2,3}{-5,4,0}{} \rgerade[0:1]{0,6,3}{0,-6,-3}{} \end{ksys3}


11g)

\begin{ksys}[xmin=-7,xmax=7,ymin=-7,ymax=7,klein] \punkt{6}{0}{} \punkt{5}{3}{} \punkt{0}{4}{} \punkt{-5}{3}{} \punkt{-6}{0}{} \punkt{-5}{-3}{} \punkt{0}{-4}{} \punkt{5}{-3}{} \end{ksys}


11h)

\begin{ksys}[xmin=-5,xmax=9,ymin=-5,ymax=9,klein] \punkt{0}{0}{A} \punkt{3}{0}{B} \punkt{3}{3}{C} \punkt{0}{3}{D} \punkt{0}{-4}{S} \punkt{-4}{0}{S} \punkt{8}{0}{S} \punkt{0}{8}{S} \end{ksys}


11i)

\begin{ksys3} \rpunkt*{3}{0}{0}{A'} \rpunkt*{5}{1}{0}{B'} \rpunkt*{1}{4}{0}{C'} \rpunkt*{-1}{3}{0}{D'} \rpunkt*{2}{2}{0}{M'} \rgerade[0:1]{3,0,0}{2,1,0}{} \rgerade[0:1]{5,1,0}{-4,3,0}{} \rgerade[0:1]{1,4,0}{-2,-1,0}{} \rgerade[0:1]{-1,3,0}{4,-3,0}{} \end{ksys3}

\erg{12}{a) Fehler: Das Minus ist verloren; es geht $4$ entgegen der $x_2$-Richtung. Richtig: $P(1 | -4 | 2)$ liegt links der $x_1x_3$-Ebene. \quad b) Die $x_2$-Koordinate ist null: Vom Ursprung aus geht man nur in $x_1$- und in $x_3$-Richtung, also bleibt $P$ in der Ebene dieser beiden Achsen.}
